\item Un estudiante sostiene un láser que emite luz de $632.8\text{ nm}$ de longitud de onda. El rayo pasa por un par de rendijas separadas $0.300\text{ mm}$, en una placa de vidrio unida al frente del láser. Después el rayo cae perpendicular en una pantalla, produciendo una configuración de interferencia. El estudiante empieza a caminar directamente hacia la pantalla a $3.00\text{ m/s}$. Si el máximo central en la pantalla permanece inmóvil, encuentre la rapidez con la que se mueven los dos máximos de primer orden ($m = 1$) en la pantalla.